\chapter{پیشینه و مبانی نظری}
\label{ch:background}

\section{پرونده الکترونیک سلامت و داده‌های مراقبت ویژه}
پرونده الکترونیک سلامت (\lr{EHR}) شامل مجموعه‌ای ناهمگون از داده‌هاست: آزمایش‌ها، علائم حیاتی، تشخیص‌ها، داروها، اقدامات، یادداشت‌ها و اطلاعات جمعیت‌شناختی. در \lr{ICU}، این داده‌ها با چگالی زمانی بالاتری ثبت می‌شوند و برای پژوهش‌های پیش‌بینی پیامد ارزش ویژه دارند. بااین‌حال، ناهمگونی ثبت، خطاهای اندازه‌گیری، مقادیر گمشده، تغییر الگوهای مراقبت و تفاوت مراکز باعث می‌شود مدل‌سازی این داده‌ها دشوار باشد.

پایگاه \lr{MIMIC\text{-}IV} یکی از مهم‌ترین منابع باز اما کنترل‌شدۀ پژوهش مراقبت ویژه است \cite{johnson2023mimic}. این پایگاه شامل داده‌های بیماران بستری در مرکز \lr{Beth Israel Deaconess Medical Center} است و از طریق \lr{PhysioNet} پس از آموزش و تعهدنامۀ استفاده از داده در دسترس قرار می‌گیرد \cite{goldberger2000physionet}. استفاده از این داده‌ها امکان بازتولید پژوهش را فراهم می‌کند، اما به‌دلیل تک‌مرکزی و گذشته‌نگر بودن، نباید با اعتبارسنجی خارجی یا آمادگی استقرار بالینی اشتباه گرفته شود.

\section{مدل‌های پیش‌بینی بالینی}
مدل پیش‌بینی بالینی باید یک جمعیت هدف، زمان پیش‌بینی، پیامد، ویژگی‌ها و معیارهای ارزیابی روشن داشته باشد. راهنماهایی مانند \lr{TRIPOD} و ابزارهایی مانند \lr{PROBAST} بر شفافیت در تعریف کوهورت، برچسب، اعتبارسنجی و خطر سوگیری تأکید می‌کنند \cite{collins2015tripod,wolff2019probast}. در این پایان‌نامه، زمان پیش‌بینی \pnum{24} ساعت نخست \lr{ICU} و پیامد، نشانگر جانشین ترخیص است.

\section{مدل‌های جدولی کلاسیک و درختی}
داده‌های ساخت‌یافته \lr{EHR} اغلب جدولی‌اند. در چنین داده‌هایی، مدل‌های کلاسیک مانند رگرسیون لجستیک و جنگل تصادفی \cite{breiman2001random} همچنان خط پایه‌های مهمی هستند. گرادیان‌بوستینگ، به‌ویژه \lr{XGBoost} \cite{chen2016xgboost}، \lr{LightGBM} \cite{ke2017lightgbm} و \lr{CatBoost} \cite{prokhorenkova2018catboost}، معمولاً روی داده‌های جدولی عملکرد قوی دارند. این مدل‌ها روابط غیرخطی و تعامل ویژگی‌ها را بدون نیاز به معماری‌های عصبی پیچیده یاد می‌گیرند.

از منظر پزشکی، تنها \lr{AUROC} کافی نیست. مدلی که احتمال‌های بدکالیبره تولید کند، حتی با تفکیک مناسب، برای تفسیر خطر قابل‌اعتماد نیست. به همین دلیل، امتیاز بریر و خطای کالیبراسیون مورد انتظار (\lr{ECE}) نیز باید بررسی شود \cite{steyerberg2019clinical,vanCalster2019calibration}.

\section{یادگیری چندنمونه‌ای}
یادگیری چندنمونه‌ای (\lr{Few-Shot Learning}) به دنبال یادگیری از تعداد اندک نمونۀ برچسب‌دار برای هر کلاس است. شبکه‌های تطبیقی و نمونه‌محور مانند \lr{Matching Networks} و \lr{Prototypical Networks} در حوزه‌هایی مانند تصویر و زبان موفق بوده‌اند \cite{vinyals2016matching,snell2017prototypical}. ایدۀ اصلی این است که مدل به‌جای یادگیری یک مرز تصمیم ثابت برای یک مجموعه بزرگ، بازنمایی‌ای بیاموزد که با چند نمونۀ پشتیبان بتواند نمونه‌های جدید را دسته‌بندی کند.

بااین‌حال، انتقال این منطق به داده‌های جدولی بالینی ساده نیست. ویژگی‌های آزمایشگاهی و علائم حیاتی معمولاً تعداد کمی بعد دارند، معنای بالینی مستقیم دارند و مدل‌های درختی روی آن‌ها قوی‌اند. بنابراین، مزیت یادگیری چندنمونه‌ای باید تجربی و در برابر خط پایه‌های قوی نشان داده شود. در این پایان‌نامه، \lr{Few-Shot ETHOS} به‌عنوان یک مسیر پژوهشی بررسی می‌شود، اما اگر ضعیف‌تر از مدل‌های جدولی باشد، این نتیجه صریحاً گزارش می‌شود.

\section{توکن‌های علائم}
توکن‌های علائم (\lr{Symptom Tokens}) تلاشی برای تبدیل داده‌های ناهمگون آزمایشگاهی و علائم حیاتی به قالبی فشرده و قابل‌پردازش هستند. هر توکن می‌تواند شناسه ویژگی، شدت، جهت تغییر یا وضعیت نرمال/غیرنرمال را رمزگذاری کند. چنین بازنمایی‌ای از نظر مفهومی به مدل‌های توکنی در پردازش زبان نزدیک است \cite{vaswani2017attention}، اما در اینجا بر داده‌های ساخت‌یافته بالینی اعمال می‌شود.

ارزش توکن‌ها ممکن است در دو سطح ظاهر شود. نخست، به‌عنوان ورودی مدل چندنمونه‌ای؛ دوم، به‌عنوان ویژگی مکمل در مدل‌های هیبریدی. نتایج این پایان‌نامه نشان می‌دهد سطح دوم در کوهورت کامل امیدوارکننده‌تر بوده است.

\section{یادگیری فدرال در سلامت}
یادگیری فدرال (\lr{Federated Learning}) اجازه می‌دهد مدل‌ها بدون انتقال مستقیم داده‌های خام میان گره‌ها آموزش ببینند \cite{mcmahan2017communication,yang2019federated}. در سلامت، این ایده با محدودیت‌های حریم خصوصی، مالکیت داده و مقررات سازگار است. بااین‌حال، فدرال بودن به‌تنهایی تضمین‌کنندۀ عملکرد یا ایمنی نیست. ناهمگونی گره‌ها، اندازۀ نمونه، ارتباطات، امنیت و اعتبارسنجی محلی باید جداگانه بررسی شوند \cite{li2019federated,bonawitz2019towards}.

در این پژوهش، یادگیری فدرال به‌صورت شبیه‌سازی‌شده انجام شد. هدف، ادعای استقرار چندمرکزی نبود، بلکه سنجش این بود که آیا یک مدل لجستیک گره‌ای می‌تواند نسبت به رگرسیون لجستیک متمرکز عملکرد قابل‌مقایسه حفظ کند یا نه.

\section{جایگاه این پایان‌نامه}
این پایان‌نامه در مرز میان انفورماتیک مراقبت ویژه، مدل‌سازی جدولی، بازنمایی توکنی، یادگیری چندنمونه‌ای و یادگیری فدرال قرار می‌گیرد. دستاورد اصلی آن یک ادعای اغراق‌آمیز درباره برتری روش نوین نیست؛ بلکه یک ارزیابی تمام‌کوهورت و قابل‌بازتولید است که نشان می‌دهد کدام روش‌ها در این وظیفه واقعاً قوی‌اند، کدام روش‌ها نیازمند بهبودند، و کدام معیارها برای تفسیر بالینی مهم‌اند.
